%% ==========================================================================
%%  FINAL REPORT --- SECTION 5, IMPLEMENTATION
%%  Sources: tech_dpf.tex, tech_serving.tex and tech_net.tex (the mid-term
%%  technical report, rewritten here because it predates private training and
%%  says so throughout), the headers under include/oblivrec/, and the output
%%  of `mingw32-make test`, which prints every number quoted below that is not
%%  from bench/results/.
%% ==========================================================================
\section{Implementation}
\label{sec:impl}

\sys{} is written in C++17 with no cryptographic library. It comprises about
7{,}400 lines of library and application code and about 3{,}300 lines of tests
in eighteen test binaries, with Python for the cleartext oracle, the export of
fixed-point models and the leakage studies. The pseudorandom generator is fixed-key
AES-128 in Davies--Meyer mode through AES-NI instructions, and key material comes
from the operating system's generator. Table~\ref{tab:layout} maps the code. The
repository also vendors \nudgeplain{}'s Go reference implementation under
\texttt{third\_party/}, under its MIT licence, as a reference to read alongside the
paper. No build target compiles or links it.

\begin{table}[t]
\centering
\small
\caption{Where each part of the system lives.}
\label{tab:layout}
\begin{tabular}{@{}l p{11.2cm}@{}}
\toprule
Directory & Contents \\
\midrule
\texttt{src/dpf/} & The DPF: tree generation, single-point and whole-domain
  evaluation, key serialisation, and the DCF built beside it. \\
\texttt{src/mpc/} & Replicated sharing, pairwise-PRF zero shares, shared
  matrix--vector products, truncation, the MSNZB comparison gate and the
  shared inverse square root. \\
\texttt{src/mf/} & Power iteration under sharing, orthogonalisation, both
  normalisers, the truncation schedule, and a cleartext fixed-point twin. \\
\texttt{src/serve/}, \texttt{src/pir/} & Score shares and the seen-item mask.
  The catalogue, the PIR server and client, and the harvest. \\
\texttt{src/net/} & Length-prefixed TCP framing, the transcript recorder, and
  channel-level network emulation. \\
\texttt{src/apps/} & The three servers, the training driver, the end-to-end
  demo, and the probe used by the distinguisher. \\
\texttt{model/} & The cleartext oracle, the float-to-ring export, the
  reconstruction attack, and the differential-privacy and poisoning studies. \\
\texttt{bench/} & Seven benchmarks writing one JSON row per measurement, and the
  script that draws every figure from those rows. \\
\bottomrule
\end{tabular}
\end{table}

\subsection{The distributed point function}
\label{sec:impl:dpf}

The DPF follows Boyle, Gilboa and Ishai~\cite{bgi2016}. It is written by us
rather than imported because it is the component both halves depend on, and a
viva will probe it hardest.

\paragraph{Key size.} A key holds a header, an initial seed, one correction word
of 18 bytes per level, and a final correction word of $\lceil \bits/8 \rceil$
bytes, so its size is $1 + 4 + 16 + 18D + \lceil \bits/8 \rceil$ bytes for a
domain of $D$ bits. The serialiser's reported size is asserted equal to the
length it actually writes. At the 11 domain bits MovieLens-100K needs, a key is
227 bytes against 16{,}384 bytes for a naive one-hot query, and the gap widens
without bound (Figure~\ref{fig:keysize}). Measuring it corrected the design
draft, which had estimated 260 bytes at 4{,}096 items. The measured value is 245
bytes, and the draft's formula had omitted the seed and the header while
evaluating at the wrong domain width.

\begin{figure}[t]
\centering
\includegraphics[width=0.62\linewidth]{fig_keysize}
\caption{Query size against catalogue size. The DPF key grows by 18 bytes per
domain bit while the naive query doubles. The star is the operating point of every
experiment in this report. The plotted closed form is the one
\texttt{tests/test\_dpf\_serialize.cpp} asserts equal to the serialised length.
Generated by \texttt{mingw32-make figures}.}
\label{fig:keysize}
\end{figure}

\paragraph{Sign convention.} Textbook presentations negate party one's output,
which gives the sum $\Eval(k_0, x) + \Eval(k_1, x) = f(x)$. Our specification
states the invariant as a difference, so we omit the negation and have
$\Eval(k_0, x) - \Eval(k_1, x) = f(x)$ at every interface, in the PIR read and in
the harvest. The DCF uses the textbook convention internally, as its construction
is defined, and converts at its boundary. Getting the convention backwards makes
every test fail in a way that looks like a traversal bug rather than a sign
error, so it is recorded at every use.

\paragraph{Key material.} The two initial seeds are the entire secret. The first
version seeded them from a Mersenne Twister, which is a genuine break: that
generator is reconstructible from 624 consecutive outputs, so an adversary who
saw enough key material could rebuild the tree and recover $\alpha$. Seeds now
come from the operating system's generator, and the failure policy is to abort
rather than fall back to a weaker source.

\paragraph{Exhaustive verification.} The characteristic failure of this
construction is a wrong control-bit update, which stays invisible until enough
levels accumulate and then presents as ``correct at two domain bits, wrong at
eight''. So the tests walk both parties' trees in lockstep and assert the
structural invariant at every internal node: on the path to $\alpha$ the seeds
differ and the control bits disagree, and off it both are equal. The difference
invariant itself was then checked for every $\alpha$ against every $x$ in every
domain up to 16 bits on the 64-bit ring and 15 bits on the 128-bit ring, which is
$7.16 \times 10^9$ leaf comparisons. That sweep passed at commit
\texttt{bce4cb1}. The default test target repeats it up to 11 bits on both rings
in about a second, and $\beta$ is derived from the domain and $\alpha$ so that
billions of comparisons do not all exercise the final correction word at one
value.

\paragraph{A bug found by testing the wire format.} Deserialisation accepted a
party byte outside $\{0, 1\}$. Evaluation starts its control bit at the party
index and branches on whether it is non-zero, so a party byte of 7 made the key
\emph{silently evaluate as party one}. Keys arrive over the network, so the party
byte and the domain width are now validated before either is used, and the tests
that found the bug are kept as regression tests.

\subsection{Private serving and delivery}
\label{sec:impl:serve}

\paragraph{The catalogue.} MovieLens-100K's item file is variable-length text, so
each record is packed into 256 bytes with explicit one-byte length prefixes
rather than delimiters. The file contains Latin-1 bytes and does not decode as
UTF-8, so it is read and carried as raw bytes. The longest record uses 117 bytes,
and every record is checked against the budget at load time, so a longer title
fails loudly at start-up rather than being cut at answer time.

\paragraph{The read.} Each PIR server expands its key over all 2{,}048 indices
and takes an inner product with the catalogue, viewed as 32 ring words per
record at $\bits = 64$. Two properties carry the obliviousness claim. The loop
visits every record in index order with no branch on the data, and every expanded
coefficient is a pseudorandom share, so there is no sparse pattern for a
compiler, a cache or a timing observer to find. Reads are checked byte for byte
against a cleartext lookup at 60 random indices on the 64-bit ring and 20 on the
128-bit ring, including both ends of the domain and padding slots.

\paragraph{Scores.} Both factors of a score are encoded at $\fracbits = 20$
fractional bits, so scores are at scale $2\fracbits$, and serving decodes at that
scale instead of truncating, because truncation is interactive. Getting this
wrong would scale every score by $2^{20}$, which is monotone and therefore
invisible in a ranking, so it is recorded explicitly. The seen-item sentinel is
$-2^{55}$, about $3{,}700\times$ below the largest real score and $256\times$
above the bottom of the ring, both asserted. The top-$\topk$ comparison is signed,
because comparing two's-complement values as unsigned would sort the masked items
first.

\paragraph{The float-to-ring boundary.} One program converts floats to ring
elements, and a committed file of 25 value pairs checks it against the C++ side
on every test run. It found two bugs on its first run. Python's \texttt{rint}
rounds half to even while C++'s \texttt{round} rounds half away from zero, so the
two disagreed at every exact half-quantum. And the reference scores had been
computed by quantising the float product, while the protocol multiplies quantised
factors, so the two disagreed on 1{,}650 of 1{,}682 scores by about $5 \times
10^{-6}$ relative. That is irrelevant to the ranking, which is exactly why a
tolerance-based check would have accepted it. Serving is now checked against the
oracle by exact equality, and all 1{,}682 scores match.

\paragraph{Three processes.} The servers are three operating-system processes
that exchange length-prefixed frames over TCP. TCP is a byte stream, so one
\texttt{send} can arrive as several \texttt{recv} calls. The test forces this with
a deliberately small receive buffer and checks that a 16{,}384-byte frame arrives
as one frame. The declared frame length is attacker-chosen and is checked against
a cap before it sizes any allocation. Every frame can be recorded at the channel
boundary into a transcript, which Section~\ref{sec:eval:wire} analyses.

\subsection{The training substrate}
\label{sec:impl:mpc}

\paragraph{Sharing and zero shares.} \texttt{Mpc3} holds all three parties' state
and moves every inter-party message through one exchange routine, which counts
rounds and bytes. That counter is what every cost in Section~\ref{sec:eval} is
read from, so a protocol cannot send without paying. Zero shares come from three
pairwise PRF generators.

\paragraph{The counter invariant.} The three generators must advance in lockstep.
A protocol that draws masks for only two of the three parties leaves the third
generator one step behind, and every later zero share then fails to sum to zero.
Nothing crashes: the parties go on producing well-formed shares of the wrong
value. This happened once and silently corrupted every subsequent multiplication.
The exchange routine now asserts that all three counters agree, and protocols in
which one party draws nothing advance its generator explicitly.

\paragraph{Theorem~4.2, asserted.} The tests check \nudgeplain{}'s central cost
claim against the byte counter rather than in prose: multiplying a shared
$4 \times 64$ matrix by a shared vector costs one round and 96 bytes, which is
three parties sending four 8-byte elements each, set by the four outputs and not
by the 256 matrix entries.

\subsection{Truncation}
\label{sec:impl:trunc}

\texttt{TruncatePair} is the standard correct construction. A helper party draws
$r$ and gives the other two parties shares of $r$ and of $r \gg \fracbits$. Those
two open $x - r$ between themselves, shift it in the clear, add back their shares
of $r \gg \fracbits$, and reshare under a pairwise mask the helper cannot
compute (Algorithm~\ref{alg:trunc}). It costs three rounds and eight ring elements
per value. Over 3{,}200 random values its error is
never more than one unit in the last place. The naive alternative, where every
party shifts its own shares, costs nothing and is catastrophically wrong on about
a quarter of values in our tests, because the shares wrap the ring. That
comparison is the whole argument for paying three rounds, and
Section~\ref{sec:eval:stages} shows the three rounds are most of what training
costs. The helper is a parameter at every call site, because a helper that also
saw the opening would learn $x$.

\subsection{The comparison gate and the normaliser}
\label{sec:impl:norm}

\paragraph{The DCF.} The design draft had asserted that a DPF and a DCF are the
same primitive. They are not, and \nudgeplain{}'s reference implementation keeps
them in separate modules. A DCF is non-zero on a whole prefix of the domain, so
every node of its tree carries a value word as well as a seed, and every level
needs a value correction word that depends on everything below it. We built it
beside the DPF, sharing the pseudorandom generator, and verified it the same
way: every threshold against every input, exhaustively up to 7 domain bits on the
64-bit ring and 6 on the 128-bit ring.

\paragraph{The MSNZB gate.} The normaliser needs the most significant non-zero
bit of a shared value $S = \norm{\mathbf{v}}^2$. The gate uses the
mask-and-reveal technique of Section~\ref{sec:background:fss}. Offline, a dealer
draws $r$ and builds one DCF key pair per threshold $2^k$, shifted by $r$. Online,
the parties open $c = S + r$ and evaluate every key at $c$. Since
$\sum_k 1[S \ge 2^k]$ counts the thresholds below $S$, any step function of
$\mathrm{msnzb}(S)$ is a public linear combination of those indicators. The seed
for $1/\norm{\mathbf{v}}$, a power of two within a factor of $\sqrt{2}$ of it, is
one such function. All
thresholds share one opening, so the gate costs two rounds, one to open and one to
reshare, however many thresholds it covers. Algorithm~\ref{alg:msnzb} gives the
gate and the Newton iteration it seeds.

\paragraph{Why the gate needs a 128-bit ring.} The opening $c = S + r$ hides $S$
only if $r$ comes from a range $\kappa$ bits wider than $S$'s. With $\kappa = 40$
and the values power iteration produces, the DCF domain needs 81 bits. A 64-bit
ring has no room for the mask at all, so \texttt{FssNormalizer} refuses to
construct at $\bits = 64$ rather than masking with too few bits, since a gate that
did so would look like privacy while providing none.

\paragraph{Inverse square root.} \texttt{InvSqrtShared} takes the gate's seed and
refines it with Newton--Raphson. Four steps reach a worst relative error of $1.41
\times 10^{-6}$ over our test values, for 50 rounds in total. The step count was
chosen from a cleartext reference first, which measured the error at every step
count before any shared code existed.

\paragraph{The gate cannot check its own range.} The gate covers thresholds from
$2^{\mathit{lo}}$ to $2^{\mathit{hi}}$. A value outside that range does not raise
an error. It clamps, because the answer is a telescoping sum over the thresholds.
And the gate cannot check, because it only ever sees $S + r$, and learning
$\mathrm{msnzb}(S)$ is precisely what it exists to avoid. So the caller must size
the range from public parameters. Section~\ref{sec:impl:bugs} describes what
happened when a caller did not.

\subsection{Power iteration}
\label{sec:impl:factor}

\texttt{ApproxFactorShared} runs Algorithm~\ref{alg:factor} with the truncation
schedule and either normaliser. On a small test instance it agrees with its
cleartext fixed-point twin to $1.5 \times 10^{-6}$ in $1 - \abs{\cos}$ per row,
with every row of $\itemembed$ unit length to $3.2 \times 10^{-6}$. The
start-up headroom assertion is not decorative. It refuses $\fracbits = 30$ at
$\bits = 64$, and it caught a real overflow: $\norm{\mathbf{v}}^2$ of a vector
not yet normalised needs 84 bits, which led to the fixed rescaling shift applied
right after the two matrix products.

The training driver writes $\itemembed$ and shares of $\userembed$ to disk, and
the demo serves them. \texttt{demo} with the privately trained, spec-faithful
model at $\bits = 128$ returns ten film titles for user 42, headed by
\emph{Toy Story}. Each record is byte-exact against a cleartext lookup, and the
ranking agrees position for position with an independently computed reference.
That run is the composition end to end: a model no server saw in the clear,
served from shares, delivered without either PIR server learning what was
fetched.

\subsection{Bugs found by measurement}
\label{sec:impl:bugs}

The project's rule was that every private component is checked against a
cleartext twin, and every approximate protocol has its error measured rather than
assumed. Table~\ref{tab:bugs} lists what that found. None of these would have
been found by reading the code, and most produced no error at all.

\begin{table}[t]
\centering
\small
\caption{Defects found by tests and measurements rather than by inspection.}
\label{tab:bugs}
\begin{tabular}{@{}p{4.3cm} p{4.4cm} p{5.8cm}@{}}
\toprule
Defect & How it was found & What it would have caused \\
\midrule
DPF party byte unvalidated & serialisation tests & a malformed key silently
  evaluates as the other party \\
Rounding mode differs across the Python and C++ boundary & 25-vector contract
  test & off-by-one ring elements at every half-quantum \\
Reference quantised the product, protocol the factors & exact score comparison &
  1{,}650 of 1{,}682 scores wrong, invisible to a tolerance check \\
PRF counters out of step & multiplication tests & well-formed shares of the wrong
  value, every product after the first \\
$\norm{\mathbf{v}}^2$ overflows before normalisation & start-up headroom
  assertion & silent wrap of the ring, a meaningless model \\
MSNZB gate range too narrow & quality study on the no-leak path & 0.108 of
  nDCG@20 lost at $\iters = 10$, with identical cost \\
Transcript writers interleaved & parsing the transcript & invalid JSON in the
  evidence file \\
\bottomrule
\end{tabular}
\end{table}

The gate-range bug deserves its own paragraph, because it is the one that nearly
became a finding. The training driver first called the MSNZB gate with the range
$[\fracbits - 8, \fracbits + 8]$, copied from the parameter sweep. Values outside
it clamped, Newton--Raphson started from the wrong power of two, and the power
iteration never converged: it ended orthogonal to the oracle's subspace. At
$\iters = 10$ the no-leak model scored 0.3428 against 0.4513 for the revealing
path, which would have been reported as ``revealing nothing costs 24\% of
recommendation quality''. The range is now $[1, 2\fracbits + 8]$ and the gap is
0.0020. Rounds and bytes are \emph{identical} under both ranges, because a wider
range buys more DCF keys offline while the online gate stays two rounds, which is
why the cost benchmark that first used the gate could never have caught it. The
rows from the buggy run were not kept, so the two nDCG values above are recorded
in the commit that fixed it (\texttt{d2138ad}) and in the comment at the gate's
call site, not in the committed data.
